%% bk-motifs.tex  表紙の図柄（各巻の数学の考え方を絵にしたもの）
%% 描く範囲：x 0〜150, y 0〜106（mm）。色 mA＝主（白か濃い色）, mB＝従（淡い色か白）, mC＝影
\makeatletter
\tikzset{mot/.style={line cap=round, line join=round}}

%% 第1巻：10のまとまり・ものさし・時計
\newcommand\bk@apple[2]{\bkapple{(#1,#2)}{4.3}{bkAppleRed}}% リンゴ（へた・葉つき）
\newcommand\bkMotifi{%
  % 10のわく：白いマスにリンゴを1こずつ
  \foreach \i in {0,...,4} \foreach \j in {0,1} {
    \fill[white, rounded corners=1.6mm] (2.6+\i*14.4,48.6+\j*14.4) rectangle ++(13.2,13.2);
    \bk@apple{9.2+\i*14.4}{54.6+\j*14.4}}
  % のこりの3こは、白いおさらの上に
  \bkplate{(101,56)}{38}
  \foreach \i in {0,1,2} \bk@apple{89+\i*12}{61.5}
  \node[mA, font=\sffamily\bfseries\fontsize{22}{22}\selectfont] at (38,39) {10};
  \node[mA, font=\sffamily\bfseries\fontsize{22}{22}\selectfont] at (101,42) {3};
  % 時計（3時）
  \fill[mB] (130,88) circle (15.5);
  \draw[mA, line width=1.6pt] (130,88) circle (15.5);
  \foreach \a in {0,30,...,330} \draw[mC, line width=1pt] (130,88) ++(\a:12.6) -- ++(\a:2.2);
  \draw[mC, line width=2pt, line cap=round] (130,88) -- ++(0:8);
  \draw[mC, line width=1.4pt, line cap=round] (130,88) -- ++(90:11.5);
  \fill[mC] (130,88) circle (1.3);
  % ものさし（0〜10）
  \fill[mB, rounded corners=1.6mm] (2,6) rectangle (122,24);
  \foreach \k in {0,...,20} {\pgfmathsetmacro\bkx{7+\k*5.5}
    \ifodd\k \draw[mC, line width=0.7pt] (\bkx,24) -- (\bkx,20.5);
    \else \draw[mC, line width=0.9pt] (\bkx,24) -- (\bkx,18.5);\fi}
  \foreach \k in {0,...,10} \node[mC, font=\sffamily\bfseries\fontsize{7.5}{8}\selectfont] at (7+\k*11,13.5) {\k};
}

%% 第2巻：落ちるボール（同じ時間ごと）・倍々・単位・数直線
\newcommand\bkMotifii{%
  \draw[mA, line width=1.4pt] (5,2) -- (5,104);
  \foreach \y in {2,12.2,...,104} \draw[mA, line width=0.9pt] (5,\y) -- (8.5,\y);
  \foreach \t/\y in {0/100, 1/93.6, 2/74.5, 3/42.7, 4/-1.9} {
    \fill[mA] (16,{\y+3}) circle (3.0);
    \node[anchor=west, mA, font=\sffamily\bfseries\fontsize{8.5}{10}\selectfont] at (20.5,{\y+3}) {\t\R{秒}{びょう}};}
  % 単位のふだ
  \foreach \x/\t in {92/m, 112/kg, 132/s} {
    \fill[mB, rounded corners=2mm] (\x,88) rectangle ++(17,13);
    \node[mC, font=\rmfamily\bfseries\fontsize{17}{17}\selectfont] at (\x+8.5,94.5) {\t};}
  % 倍々のブロック（1, 2, 4, 8）
  \begin{scope}[shift={(41,44)}]
    \foreach \bx/\bw/\bh/\bn in {0/1/1/1, 19.5/2/1/2, 47.5/2/2/4, 75.5/4/2/8} {
      \fill[mA] (\bx,0) rectangle ++(\bw*8.5,\bh*8.5);
      \draw[mC, line width=0.7pt, shift={(\bx,0)}] (0,0) grid[step=8.5] ({\bw*8.5},{\bh*8.5});
      \node[mA, font=\sffamily\bfseries\fontsize{16}{16}\selectfont] at (\bx+\bw*4.25,\bh*8.5+6.5) {\bn};}
    \foreach \xm in {14, 42, 70} \node[mB, font=\sffamily\bfseries\fontsize{11}{11}\selectfont] at (\xm,4.25) {×2};
  \end{scope}
  % 数直線（−5〜5）
  \draw[mA, line width=1.6pt, -{Stealth[length=8pt]}] (40,14) -- (150,14);
  \foreach \k in {-5,...,5} {\pgfmathsetmacro\bkx{93+\k*10}
    \draw[mA, line width=1.1pt] (\bkx,11.5) -- (\bkx,16.5);
    \node[mA, font=\sffamily\bfseries\fontsize{9.5}{9.5}\selectfont] at (\bkx,6.5) {\ifnum\k<0 −\number-\k\else\k\fi};}% 第7回：数字をゴシックに（マイナスは記号の字）
  \draw[mB, line width=3.4pt, -{Stealth[length=9pt]}] (93,24) -- (123,24);
  \draw[mB, line width=3.4pt, -{Stealth[length=9pt]}] (93,24) -- (73,24);
  \fill[mA] (93,24) circle (2.4);
}

%% 第3巻：データに直線・投げたボールの道・ヒストグラム
\newcommand\bkMotifiii{%
  \draw[mA, line width=1.3pt, -{Stealth[length=7pt]}] (4,6) -- (98,6);
  \draw[mA, line width=1.3pt, -{Stealth[length=7pt]}] (6,4) -- (6,104);
  \draw[mB, line width=4pt] (10,12) -- (94,86);
  \foreach \x/\e in {15/3, 24/-4, 33/2.5, 42/-2, 51/4, 60/-3, 69/1.5, 78/-1.5, 87/2.5}
    \fill[mA] (\x,{12+(\x-10)*74/84+\e}) circle (2.3);
  \draw[mA, line width=1.3pt, dash pattern=on 4pt off 3pt] plot[domain=104:148, samples=40] (\x,{8+74*(1-((\x-126)/22)^2)});
  \foreach \x in {104,111.33,118.67,126,133.33,140.67,148} \fill[mA] (\x,{8+74*(1-((\x-126)/22)^2)}) circle (2.4);
  \foreach \h [count=\i] in {2,5,10,15,18,15,10,5,2}
    \fill[mB] (99+\i*5.4,88) rectangle ++(4.8,\h);
  \draw[mA, line width=1pt] (103,88) -- (150,88);
}

%% 第4巻：円運動の影＝波・ベクトルの合成
\newcommand\bkMotifiv{%
  \def\bkcy{64}
  \draw[mB, line width=0.9pt] (2,\bkcy) -- (66,\bkcy) (34,\bkcy-32) -- (34,\bkcy+32);
  \draw[mA, line width=1.6pt] (34,\bkcy) circle (28);
  \fill[mB] (34,\bkcy) -- (42,\bkcy) arc (0:50:8) -- cycle;
  \draw[mA, line width=2pt] (34,\bkcy) -- ++(50:28);
  \draw[mC, line width=1.4pt, dash pattern=on 3pt off 2pt] (52,{\bkcy+28*sin(50)}) -- (52,\bkcy);
  \draw[mC, line width=3pt] (34,\bkcy) -- (52,\bkcy);
  \fill[mA] (52,{\bkcy+28*sin(50)}) circle (2.4);
  \draw[mA, line width=1pt, dash pattern=on 3pt off 2pt] (54,{\bkcy+28*sin(50)}) -- (72,{\bkcy+28*sin(50)});
  \draw[mB, line width=0.9pt] (72,\bkcy) -- (150,\bkcy);
  \draw[mA, line width=2.2pt] plot[domain=72:150, samples=120, smooth] (\x,{\bkcy+28*sin(50+(\x-72)*5.142857)});
  % ベクトル
  \begin{scope}[shift={(92,4)}]
    \draw[mB, line width=1pt, dash pattern=on 3pt off 2pt] (32,7) -- (44,22) -- (12,15);
    \draw[mA, line width=2pt, -{Stealth[length=8pt]}] (0,0) -- (32,7);
    \draw[mA, line width=2pt, -{Stealth[length=8pt]}] (0,0) -- (12,15);
    \draw[mC, line width=2.6pt, -{Stealth[length=9pt]}] (0,0) -- (44,22);
  \end{scope}
}

%% 第5巻：接線（微分）と細い長方形（積分）
\newcommand\bkMotifv{%
  \draw[mA, line width=1.3pt, -{Stealth[length=7pt]}] (4,6) -- (150,6);
  \draw[mA, line width=1.3pt, -{Stealth[length=7pt]}] (6,4) -- (6,104);
  \foreach \k in {0,...,7} {\pgfmathsetmacro\bkx{16+\k*10}
    \pgfmathsetmacro\bkh{16+0.35*((\bkx-6)/10)^2+9*sin(3*(\bkx-6))}
    \fill[mB] (\bkx,6) rectangle (\bkx+10,\bkh);
    \draw[mC, line width=0.6pt] (\bkx,6) rectangle (\bkx+10,\bkh);}
  \draw[mA, line width=2.4pt] plot[domain=8:146, samples=120, smooth] (\x,{16+0.35*((\x-6)/10)^2+9*sin(3*(\x-6))});
  \pgfmathsetmacro\bky{16+0.35*11.4^2+9*sin(3*114)}
  \pgfmathsetmacro\bks{0.007*114+9*3*pi/180*cos(3*114)}
  \draw[mA, line width=1.2pt, dash pattern=on 3pt off 2pt] (120,\bky) -- (134,\bky) -- (134,\bky+14*\bks);% 第6回：緑地で見えるように白に
  \draw[mA, line width=1.8pt] (100,{\bky-20*\bks}) -- (146,{\bky+26*\bks});
  \fill[mA] (120,\bky) circle (2.6);
  \node[mA, anchor=north, font=\sffamily\bfseries\fontsize{10}{10}\selectfont] at (127,\bky-1.8) {$\Delta x$};% 第6回：横の辺の中央の下
  \node[mA, anchor=west, font=\sffamily\bfseries\fontsize{10}{10}\selectfont] at (135,\bky+8) {$\Delta y$};
}

%% 第6巻：ゆれながら止まる（減衰振動）・回る矢印（複素数）
\newcommand\bkMotifvi{%
  \draw[mB, line width=0.9pt] (6,56) -- (150,56);
  \draw[mA, line width=1.3pt, -{Stealth[length=7pt]}] (6,6) -- (6,106);
  \draw[mB, line width=1.2pt, dash pattern=on 4pt off 3pt] plot[domain=8:148, samples=80] (\x,{56+40*exp(-0.022*(\x-8))});
  \draw[mB, line width=1.2pt, dash pattern=on 4pt off 3pt] plot[domain=8:148, samples=80] (\x,{56-40*exp(-0.022*(\x-8))});
  \draw[mA, line width=2.4pt] plot[domain=8:148, samples=200, smooth] (\x,{56+40*exp(-0.022*(\x-8))*cos(9.17*(\x-8))});
  \begin{scope}[shift={(132,90)}]
    \draw[mB, line width=0.8pt] (-15,0) -- (15,0) (0,-14) -- (0,15);
    \draw[mA, line width=1.2pt] (0,0) circle (11);
    \draw[mA, line width=2pt, -{Stealth[length=7pt]}] (0,0) -- (40:11);
    \draw[mB, line width=1.4pt, -{Stealth[length=6pt]}] (55:13.5) arc (55:120:13.5);
  \end{scope}
}

%% 第7巻：行列で変わる格子
\newcommand\bkMotifvii{%
  \begin{scope}
    \clip (0,0) rectangle (150,106);
    \begin{scope}[shift={(72,50)}]
      \foreach \i in {-7,...,7} {
        \draw[mB, line width=0.5pt] (\i*12,-60) -- (\i*12,60);
        \draw[mB, line width=0.5pt] (-90,\i*12) -- (90,\i*12);}
      \begin{scope}[cm={1,0.28,0.48,1,(0,0)}]
        \foreach \i in {-8,...,8} {
          \draw[mA, line width=0.9pt] (\i*12,-80) -- (\i*12,80);
          \draw[mA, line width=0.9pt] (-110,\i*12) -- (110,\i*12);}
        \fill[mC] (0,0) rectangle (12,12);
        \draw[mA, line width=2.4pt, -{Stealth[length=9pt]}] (0,0) -- (24,0);
        \draw[mA, line width=2.4pt, -{Stealth[length=9pt]}] (0,0) -- (0,24);
      \end{scope}
      \fill[mA] (0,0) circle (2);
    \end{scope}
  \end{scope}
}

%% 第8巻：わき出し（div）と うず（rot）
\newcommand\bkMotifviii{%
  % うすい等高線と、いちばん急な向き（grad）
  \foreach \r in {6,12,...,30} \draw[mB, line width=0.5pt] (75,94) ellipse ({\r*1.5} and {\r*0.3});
  % わき出し
  \begin{scope}[shift={(36,50)}]
    \fill[mB] (0,0) circle (5);\fill[mA] (0,0) circle (2.4);
    \foreach \a in {0,30,...,330} \draw[mA, line width=1.6pt, -{Stealth[length=5pt]}] (\a:9) -- (\a:16);
    \foreach \a in {15,45,...,345} \draw[mA, line width=1.3pt, -{Stealth[length=5pt]}] (\a:21) -- (\a:28);
    \node[mB, font=\sffamily\bfseries\fontsize{14}{14}\selectfont] at (0,-38) {div};
  \end{scope}
  % うず
  \begin{scope}[shift={(112,50)}]
    \fill[mB] (0,0) circle (2.4);
    \foreach \a in {0,45,...,315} \draw[mA, line width=1.6pt, -{Stealth[length=5pt]}] (\a:11) ++(\a-90:2.6) -- ++(\a+90:5.2);
    \foreach \a in {15,45,...,345} \draw[mA, line width=1.3pt, -{Stealth[length=5pt]}] (\a:24.5) ++(\a-90:3.6) -- ++(\a+90:7.2);
    \node[mB, font=\sffamily\bfseries\fontsize{14}{14}\selectfont] at (0,-38) {rot};
  \end{scope}
}

%% 第9巻：sinをたして四角い波・スペクトル
\newcommand\bkMotifix{%
  \def\bkmid{72}\def\bkamp{15}
  \draw[mB, line width=1.2pt, dash pattern=on 4pt off 3pt] (8,\bkmid-\bkamp) -- (8,\bkmid+\bkamp) -- (43,\bkmid+\bkamp) -- (43,\bkmid-\bkamp) -- (78,\bkmid-\bkamp) -- (78,\bkmid+\bkamp) -- (113,\bkmid+\bkamp) -- (113,\bkmid-\bkamp) -- (148,\bkmid-\bkamp) -- (148,\bkmid);
  \draw[mB, line width=1.1pt] plot[domain=8:148, samples=120, smooth] (\x,{\bkmid+\bkamp*4/pi*sin((\x-8)*5.142857)});
  \draw[mA, line width=2.2pt] plot[domain=8:148, samples=320, smooth] (\x,{\bkmid+\bkamp*4/pi*(sin((\x-8)*5.142857)+sin(3*(\x-8)*5.142857)/3+sin(5*(\x-8)*5.142857)/5+sin(7*(\x-8)*5.142857)/7+sin(9*(\x-8)*5.142857)/9)});
  \draw[mA, line width=1.3pt] (8,8) -- (148,8);
  \foreach \k [count=\i] in {1,3,5,7,9,11} {
    \fill[mA] ({4+\i*18},8) rectangle ++(9,{36/\k});
    \node[mB, font=\sffamily\bfseries\fontsize{8}{8}\selectfont] at ({8.5+\i*18},4) {\k};}
}
\makeatother
